\section{Задача и сетевые участники}
\begin{definition}[Вариант 34: сортировочный центр]
Посылки с уникальными номерами проходят входной накопитель, сканирование, сортировку, конвейер и выход своего направления. Вместимость очередей и число устройств ограничены. При заполнении следующего накопителя устройство удерживает посылку. Ошибки маркировки вызывают повторное чтение, после исчерпания попыток - передачу в ручной накопитель.
\end{definition}

Программы написаны на C11 для Linux. Сервер хранит модель и выдаёт задания по TCP. Рабочие клиенты моделируют источники, сканеры, сортировщики, конвейеры и разгрузчики. Каждый клиент - отдельный процесс; посылки представлены записями сервера с собственной стадией и состоянием маркировки. Общей памяти между приложениями нет.

По умолчанию работают 13 рабочих клиентов, сервер и наблюдатели. В каждом приложении один цикл \texttt{poll}, без потоков pthread.

\section{Реализованные уровни требований}
\begin{center}
\begin{tabular}{>{\raggedright\arraybackslash}p{0.17\linewidth}>{\raggedright\arraybackslash}p{0.73\linewidth}}
\toprule
Уровень & Реализация \\
\midrule
4--5 & Независимые TCP-программы, IP и порт в аргументах, события своего обмена, запуск на разных компьютерах \\
6--7 & Отдельный наблюдатель: полный снимок и общий поток событий \\
8 & Несколько наблюдателей, отключение и новое подключение без остановки модели \\
Бонусы & Возврат рабочих клиентов с сохранением задания; завершение подключённых клиентов при остановке сервера \\
\bottomrule
\end{tabular}
\end{center}
Режимы отличаются числом наблюдателей. Подробности - \path{docs/PROTOCOL.md}.

\section{Обмен, корректность и восстановление}
Протокол использует строки с LF, накопление неполных сообщений и буферы частичной отправки. Переполнение исходящего буфера отключает только медленного клиента.

\begin{remark}[Сохранение задания]
После \texttt{HELLO} сервер связывает клиента с ролью и номером. \texttt{JOB} содержит токен, посылку, длительность и результат моделируемой операции. Клиент выдерживает задержку и посылает \texttt{DONE}. Один токен изменяет состояние только один раз. При разрыве незавершённое задание сохраняется; после возврата отправляется снова, а его задержка начинается заново. Уже принятый результат повторно не применяется.
\end{remark}

Очереди FIFO. Повторное чтение остаётся на том же сканере, чтобы не возвращать посылку в переполненный вход. Проверяются единственное расположение посылки, постоянство ID, вместимости, занятость устройств, распознавание перед выходом и итоговый баланс.

Первоначальный старт ожидает все рабочие роли. Наблюдатель получает снимок между \texttt{BEGIN}/\texttt{END} с одинаковым номером события, затем новые события. Heartbeat раз в секунду обнаруживает молчащее соединение за 5--6 секунд.

\newpage
\section{Запуск, интерфейс и завершение}
\begin{lstlisting}
make
./build/center_server 127.0.0.1 9034 \
    --delay-ms 30 --check --log events.txt
# in separate terminals:
./build/center_observer 127.0.0.1 9034
python3 scripts/run_workers.py 127.0.0.1 9034
\end{lstlisting}
Клиент: \texttt{center\_worker IP PORT scanner 1}; набор ролей указан в \texttt{PLAN}. Для нескольких машин сервер слушает \texttt{0.0.0.0}, клиенты используют его доступный IPv4. Python автоматизирует запуск C-программ и тесты.

Параметры перечислены в \texttt{-{}-help}. Пределы: 512 посылок, 16 устройств каждого вида, 256 соединений. Seed фиксирует результаты посылок, но не расписание.

Вывод показывает время, ID, роль, стадию, очередь, чтение, повторы, ожидания и итоги. Наблюдатель дополнительно получает полное состояние модели; \texttt{-{}-log} сохраняет серверный журнал.

После доставки или отказа всей партии сервер отправляет итог и \texttt{STOP complete}. SIGINT/SIGTERM отправляют \texttt{STOP signal} и частичную статистику. Сервер завершает передачу через \texttt{shutdown}, дочитывает ответы не более двух секунд и закрывает сокеты. Клиенты завершаются по STOP. Ctrl+C отдельного рабочего клиента отключает только его.

Простой ограничен \texttt{-{}-startup-timeout} (30 с), повторы подключения - \texttt{-{}-retry-seconds} (10 с). Сервер возвращает 0 при успехе, 130/143 при сигнале, 1 при ошибке, 2 при неверных аргументах. Клиенты после STOP возвращают 0.

\section{Проверки и результаты}
Linux ARM64 в Docker, GCC 12.2, строгая сборка \texttt{-Wall -Wextra -Wpedantic -Werror}. Тесты независимо восстанавливают FIFO, владельцев посылок и допустимые переходы по событиям.
\begin{center}
\begin{tabular}{>{\raggedright\arraybackslash}p{0.58\linewidth}>{\raggedright\arraybackslash}p{0.32\linewidth}}
\toprule
Проверка & Результат \\
\midrule
Полный набор, включая сетевые отказы & 30 из 30, 39.010 с \\
Тот же набор с ASan и UBSan & 30 из 30, 108.008 с \\
40 посылок, без ошибок маркировки & 40 доставок, 40 чтений \\
17 посылок, 100\% ошибок, 3 повтора & 17 отказов, 68 чтений \\
Три контейнера с разными сетевыми пространствами & 30 доставок, все коды 0 \\
\bottomrule
\end{tabular}
\end{center}
Проверены 512 посылок, случайные конфигурации, медленный наблюдатель, дубликаты, частичные сообщения, возврат каждой роли и сигналы. Вывод - \path{docs/test-results.txt}. Физические машины отдельно не использовались; проверено разделение контейнерами.

\section{Критическая оценка и исправления}
Исправлены раннее завершение короткой партии (барьер регистрации), потеря STOP при \texttt{POLLHUP} (дочитывание сокета) и бесконечные повторы занятой роли (ограничение срока).

Очереди - $O(1)$, проверка - $O(N+J+Q)$, снимок - $O(N(N+J+Q))$. Буферы занимают около 32 МиБ. Сервер - единая точка отказа: после SIGKILL клиенты обнаруживают разрыв, но модель не восстанавливается. Сеть считается доверенной, TLS нет.

Из ИДЗ3 сохранены параметры, алгоритм и оформление. ИИ использован для кода, тестов и отчёта: \path{docs/PROMPTS.md}. Исправления: \path{docs/DEBUGGING.md}.
